\chapter[\emj{🧬}~LCA y Serialización de Árboles]{\emj{🧬}~Ancestro Común (LCA) y Serialización de Árboles}

\begin{dsaquees}

Dos primos buscando a su pariente compartido más cercano en un árbol
genealógico 🌳. El Ancestro Común más Bajo (Lowest Common Ancestor) es
la primera persona de la que ambos descienden: ni muy arriba (el
bisabuelo de todos), ni imposible (uno mismo).

Y la Serialización es tomarle una "foto plana" al árbol para guardarlo
en un archivo o mandarlo por la red, de forma que después puedas
RECONSTRUIRLO idéntico. Como desarmar un mueble con instrucciones para
volverlo a armar exactamente igual.

\end{dsaquees}

\begin{dsaotraforma}

El \textbf{LCA} de dos nodos $p$ y $q$ es el nodo más profundo que tiene
a ambos como descendientes (un nodo puede ser descendiente de sí mismo).

\begin{itemize}
\item En un \textbf{BST}: aprovecha el orden. Baja desde la raíz; si
      ambos valores son menores ve a la izquierda, si ambos son mayores
      ve a la derecha; en cuanto se "separan", ese nodo es el LCA.
\item En un \textbf{árbol binario general}: DFS recursivo. Si encuentras
      $p$ o $q$ los devuelves; el primer nodo que recibe respuesta
      distinta de NULL por AMBOS lados es el LCA.
\end{itemize}

La \textbf{serialización} convierte el árbol en una cadena (típicamente
un preorder con marcadores para los NULL) y la deserialización la
reconstruye. Es la base de guardar árboles en disco o red.

\end{dsaotraforma}

\begin{dsadiagrama}
Árbol:
           [3]
         /     \
       [5]     [1]
      /   \    /  \
    [6]  [2] [0]  [8]

LCA(6, 2) = 5     (ambos cuelgan de 5)
LCA(6, 8) = 3     (se separan en la raíz)
LCA(5, 2) = 5     (un nodo es ancestro de sí mismo)

Serialización (preorder, # = NULL):
  3,5,6,#,#,2,#,#,1,0,#,#,8,#,#
Deserializar lee en el mismo orden y reconstruye idéntico.
\end{dsadiagrama}

\begin{dsacomofunciona}
\begin{verbatim}
LCA(6, 8) subiendo la señal desde las hojas

           (3) ◄── ambos lados devuelven algo → LCA = 3
          /   \
    izq (5)    (1) der
       / \     /  \
    (6)  2   0   (8)
     ▲            ▲
     p            q

dfs(3): izq=dfs(5)  der=dfs(1)
  dfs(5): izq=dfs(6)=6 OK   der=dfs(2)=null   → devuelve 6
  dfs(1): izq=dfs(0)=null  der=dfs(8)=8 OK    → devuelve 8
  en 3: izq=6 (≠null) Y der=8 (≠null) → 3 es el LCA OK

SERIALIZAR (preorder, # = NULL) produce un stream lineal:
  3 , 5 , 6 , # , # , 2 , # , # , 1 , 0 , # , # , 8 , # , #
  └─raíz┘ └──subárbol izq──┘        └────subárbol der────┘
Deserializar consume el stream en el MISMO orden → árbol idéntico.
\end{verbatim}
\end{dsacomofunciona}

\begin{lstlisting}
LCA EN ÁRBOL BINARIO GENERAL (DFS)
1. Caso base: si el nodo es NULL o es p o q, devuélvelo.
2. Busca en izquierda y en derecha.
3. Si ambos lados devuelven algo -> este nodo es el LCA.
4. Si solo un lado devuelve algo -> propaga ese resultado hacia arriba.

PSEUDOCÓDIGO (LCA general)
──────────────────────────────────────────────────────────────

función lca(nodo, p, q):
    SI nodo == NULL O nodo == p O nodo == q: return nodo
    izq = lca(nodo.izq, p, q)
    der = lca(nodo.der, p, q)
    SI izq != NULL Y der != NULL: return nodo   // p y q en lados distintos
    return (izq != NULL) ? izq : der

CÓDIGO C++ (LCA general y en BST)
──────────────────────────────────────────────────────────────

struct Node { int val; Node *left, *right; };

// Árbol binario general -> O(n)
Node* lca(Node* root, Node* p, Node* q) {
    if (!root || root == p || root == q) return root;
    Node* l = lca(root->left,  p, q);
    Node* r = lca(root->right, p, q);
    if (l && r) return root;           // se separan aquí -> LCA
    return l ? l : r;
}

// BST -> O(h), aprovecha el orden
Node* lcaBST(Node* root, int p, int q) {
    while (root) {
        if (p < root->val && q < root->val)      root = root->left;
        else if (p > root->val && q > root->val) root = root->right;
        else return root;              // se separan -> LCA
    }
    return nullptr;
}

CÓDIGO C++ (serializar / deserializar, preorder)
──────────────────────────────────────────────────────────────

#include <string>
#include <sstream>
using namespace std;

void serializeH(Node* n, string& out) {
    if (!n) { out += "#,"; return; }
    out += to_string(n->val) + ",";
    serializeH(n->left, out);
    serializeH(n->right, out);
}

Node* deserializeH(istringstream& in) {
    string tok; getline(in, tok, ',');
    if (tok == "#") return nullptr;
    Node* n = new Node{stoi(tok), nullptr, nullptr};
    n->left  = deserializeH(in);
    n->right = deserializeH(in);
    return n;
}

COMPLEJIDAD (Big-O)
──────────────────────────────────────────────────────────────

  Operación         │ Tiempo  │ Espacio
  ──────────────────┼─────────┼──────────
  LCA (general)     │ O(n)    │ O(h)
  LCA (BST)         │ O(h)    │ O(1)
  Serializar        │ O(n)    │ O(n)
  Deserializar      │ O(n)    │ O(n)
\end{lstlisting}

\begin{dsaentrevistas}

✅ "Lowest Common Ancestor" (LeetCode 235 BST / 236 general): dos
   versiones. En BST usa el orden (más rápido, O(h)); en árbol general
   usa el DFS de "ambos lados devuelven algo".

💡 "Serialize and Deserialize Binary Tree" (LeetCode 297): clave usar un
   marcador para NULL (\verb|#|) para que la reconstrucción sea única.

⚠️ Sin marcador de NULL no se puede reconstruir: un preorder de solo
   valores es ambiguo. Los \verb|#| son obligatorios.

🔑 Distancia entre dos nodos = profundidad(p) + profundidad(q) −
   2·profundidad(LCA). El LCA es la pieza para medir distancias en
   árboles.

\end{dsaentrevistas}

\begin{dsaresumen}

• LCA: ancestro común más profundo de dos nodos.
• BST: usa el orden, baja hasta que p y q se separan; O(h).
• General: DFS; el nodo donde ambos lados devuelven algo es el LCA.
• Serializar: preorder con \verb|#| para los NULL; O(n).
• Deserializar lee en el mismo orden y reconstruye idéntico.
• LCA es la base para medir distancias entre nodos.
════════════════════════════════════════════════════════════════

\end{dsaresumen}
